% !TeX spellcheck = es_ES
\documentclass[12pt]{article}

%\usepackage[utf8]{inputenc}
\usepackage{amsmath}
\usepackage{moodle}
\begin{document}
\begin{quiz}{Matemáticas/Límites}
% \begin{multi}[numbering=none, 
% feedback={
% \[x^2+1\]
% y hacemos Taylor
% }]{A first derivative}
%   What is the first derivative of $x^3$?
%   \item  $\frac{1}{4} x^4+C$
%   \item* $3x^2$
%   \item  $51$
% \end{multi}
 
 \begin{multi}[shuffle=true,numbering=none,
 feedback={Los límites laterales de la función en cero no coinciden. Por tanto, la función no es continua en cero.}]{Pregunta continuidad}
	La función 
	\[
	f(x)= \begin{cases}
		x^2-x, & \text{si $x <0$}, \\
		1+\log(x+1), &  \text{si $x\geq 0$},
		\end{cases}
	\]
	es continua.
	\item Verdadero 
	\item* Falso
 \end{multi}
 
  \begin{multi}[shuffle=true,numbering=none,
  feedback={Dividiendo en el numerador y en el denominador por la máxima potencia, que es $x^2$, nos queda el límite siguiente:
  \[
  \lim_{x \to +\infty}\frac{2-\frac{5}{x}+\frac{14}{x^2}}{4-\frac{5}{x}+\frac{19}{x^2}}=\frac{2}{4}= \frac{1}{2}
  \]
  }]{Pregunta límites}Calcula el siguiente límite: $\displaystyle \lim_{x \to +\infty} \frac{2x^2-5x+14}{4x^2-5x+19}$.
  
  \item $+\infty$ 
  \item* $\frac{1}{2}$  
  \item $2$
  \end{multi}
  
  \begin{multi}[shuffle=true,numbering=none,
    feedback={El límite propuesto presenta una indeterminación de ``$\frac{0}{0}$''; factorizando la expresión nos queda:
    \[
    \frac{x(x-1)^2}{x(x-1)(x+1)}=\frac{\cancel{x}\cancel{(x-1)}(x-1)}{\cancel{x}\cancel{(x-1)}(x+1)}=\frac{x-1}{x+1}
    \]
    así que cuando $x \to 1$, la función tiende a $0$.
    }]{Pregunta límites}Calcula $\displaystyle \lim_{x\to 1}\frac{x^3-2x^2+x}{x^3-x}$.
 
	 \item* $0$ 
	 \item $1$ 
	 \item $+\infty$
	\end{multi}
	
	 \begin{multi}[shuffle=true,numbering=none,
	    feedback={Dividiendo en el numerador y en el denominador por la máxima potencia, que es $x^3$, nos queda el límite siguiente:
	    \[
	    \lim_{x \to -\infty}\frac{\frac{4}{x}-\frac{3}{x^2}+\frac{7}{x^3}}{1-\frac{2}{x}+\frac{5}{x^2}}=\frac{0}{1}= 0.
	    \]
	    }]{Pregunta límites}Calcula $\displaystyle \lim_{x \to -\infty}\frac{4x^2-3x+7}{x^3-2x^2+5x}$.
	    \item* $0$ 
	    \item $+\infty$  
	    \item $-\infty$
	   \end{multi}
	   
	   \begin{multi}[shuffle=true,numbering=none,
	   	    feedback={Multiplicamos el numerador y el denominador por el conjugado del numerador ($\sqrt{x}+\sqrt{x+1}$) y nos queda el límite siguiente:
	   	    \[
	   	    \lim_{x \to +\infty}\frac{x-(x+1)}{x^2(\sqrt{x}+\sqrt{x+1})}=\lim_{x \to +\infty}\frac{-1}{x^2(\sqrt{x}+\sqrt{x+1})}= 0.
	   	    \]
	   	    }]{Pregunta límites}Calcula el siguiente límite: $\displaystyle \lim_{x \to +\infty} \frac{\sqrt{x}-\sqrt{x+1}}{x^2}$.
			\item* $0$ 
			\item $\frac{1}{4}$   
			\item $-\infty$
	   \end{multi}
	   
	   \begin{multi}[shuffle=true,numbering=none,
	   	   	    feedback={El límite presenta una indeterminación de ``$\frac{0}{0}$'' por lo que vamos a factorizar:
	   	   	    \[
	   	   	    \lim_{x \to 1}\frac{\sqrt{x-1}}{\sqrt{x^2-x}}=\lim_{x \to 1} \sqrt{\frac{x-1}{x(x-1)}}=
	   	   	    \lim_{x \to 1} \sqrt{\frac{\cancel{(x-1)}}{x\cancel{(x-1)}}} =\lim_{x \to 1} \sqrt{\frac{1}{x}}= 1.
	   	   	    \]
	   	   	    }]{Pregunta límites}Calcula el siguiente límite: $\displaystyle \lim_{x \to 1} \frac{\sqrt{x-1}}{\sqrt{x^2-x}}$.
	   	   
		   \item $+\infty$ 
		   \item $0$   
		   \item* $1$
	   \end{multi}
	   
 	   \begin{multi}[shuffle=true,numbering=none,
  	   	    feedback={El límite no presenta ninguna una indeterminación ya que el primer sumando tiende a $+\infty$ y el segundo sumando a 1. Por tanto, el resultado es:
  	   	    \[
  	   	    \lim_{x \to 0^+}\left(\frac{1}{x}-\frac{1}{x^2+1}\right)=+\infty.
  	   	    \]
  	   	    }]{Pregunta límites}Calcula el siguiente límite: $\displaystyle \lim_{x \to 0^+} \left(\frac{1}{x}-\frac{1}{x^2+1}\right)$.
	   	   	   	
	   	  \item $0$ 
	   	  \item* $+\infty$   
	   	  \item $-\infty$
		 \end{multi}
		 
		 \begin{multi}[shuffle=true,numbering=none,
		   	   	    feedback={El límite  presenta  una indeterminación del tipo ``$\infty-\infty$''ya que el primer sumando tiende a $+\infty$ y el segundo sumando a $+\infty$. Tomamos común denominador y multiplicando el numerador y denominador por $\sqrt{x}$ el resultado es:
		   	   	    \[\begin{aligned}
		   	   	    \lim_{x \to 0^+}\left(\frac{1}{x}-\frac{1}{\sqrt{x}}\right)&=\lim_{x \to 0^+} \frac{\sqrt{x}-x}{x\sqrt{x}}=\lim_{x \to 0^+} \frac{x-x\sqrt{x}}{x^2} \\
		   	   	    &= \lim_{x \to 0^+} \frac{x(1-\sqrt{x})}{x^2}= \lim_{x \to 0^+} \frac{1-\sqrt{x}}{x}=+\infty.
		   	   	    \end{aligned}\]
		   	   	    }]{Pregunta límites}Calcula el siguiente límite: $\displaystyle \lim_{x \to 0^+} \left(\frac{1}{x}-\frac{1}{\sqrt{x}}\right)$.
		   	   	    
		   	\item* $+\infty$ 
		   	\item $-\infty$   
		   	\item $0$
		  \end{multi}
		  
		  
		  \begin{multi}[shuffle=true,numbering=none,
		  		   	   	    feedback={\[
		  		   	   	    lim_{x \to 0^+} \frac{\sqrt{x}}{x} = \lim_{x \to 0^+} \sqrt{\frac{x}{x^2}} = \lim_{x \to 0^+} \sqrt{\frac{1}{x}} = +\infty
		  		   	   	    \]
		  		   	   	    }]{Pregunta límites}Calcula el siguiente límite: $\displaystyle \lim_{x \to 0^+} \frac{\sqrt{x}}{x}$.
			  \item  $-1$ 
			  \item $1$ 
			  \item* $+\infty$
		  \end{multi}
		  
		  \begin{multi}[shuffle=true,numbering=none,
		  		  		   	   	feedback={Como $\lim_{x \to 0} f(x)=+\infty$, la función tiene un asíntota vertical en $x=0$. La recta $y=1$ es una asíntota horizontal ya que $\lim_{x \to \infty} f(x)=1$. No tiene asíntotas oblicuas.
		  		  		   	   	 }]{Asíntotas}La función $f(x)=x+\frac{1}{x^2}$,
		  
			  \item*  tiene una asíntota vertical en $x=0$ 
			  \item tiene una asíntota horizontal en $y=0$ 
			  \item no tiene asíntotas
		  \end{multi}
		  
		  \begin{multi}[shuffle=true,numbering=none,
		  		  		  		   	   	feedback={Como $\lim_{x \to \infty} f(x) =1/3$, la recta $y=1/3$ es una asíntota horizontal.
		  		  		  		   	   	}]{Asíntota vertical}La función $f(x)=\frac{x^2+x-7}{3x^2+x+1}$,
			  \item $x=1/3$ es una asíntota vertical 
			  \item* $y=1/3$ es una asíntota horizontal
		  \end{multi}
		  
		  \begin{multi}[shuffle=true,numbering=none,
		  		  		  		  		   	   	feedback={}]{Asíntota oblícua}La función $f(x)=\frac{4x^4-2x^2+x}{x^3-2x+1}$ tiene una asíntota oblicua $y=mx+n$. 
			  	\item $m=4$, $n=1$ 
			  	\item* $m=4$, $n=0$ 
			  	\item $m=\frac{1}{4}$, $n=0$	  		  		  		   	   			  
		  \end{multi}
 \end{quiz}
\end{document}

